\documentclass{article}

\usepackage[utf8]{inputenc}
\usepackage[T1]{fontenc}
\usepackage[english]{babel}
\usepackage{microtype}

\usepackage{amsmath, amssymb, amsthm}
\usepackage{mathtools}
\usepackage{totcount}
\usepackage{enumitem}

\usepackage{newtxtext,newtxmath}

\usepackage{titlesec}
\titleformat{\chapter}[display]
  {\normalfont\huge\bfseries} % Format of the text
  {Lecture \thechapter}       % Label (e.g., "Lecture 1")
  {0.5em}                     % Spacing between label and title
  {\Huge}                     % Format of the title

% Theorem-like environments
\newtheorem{theorem}{Theorem}[section]
\newtheorem{lemma}[theorem]{Lemma}
\newtheorem{proposition}[theorem]{Proposition}
\newtheorem{corollary}[theorem]{Corollary}

\theoremstyle{definition}
\newtheorem{definition}[theorem]{Definition}
\newtheorem{example}[theorem]{Example}

\theoremstyle{remark}
\newtheorem{remark}[theorem]{Remark}

\DeclareMathOperator*{\lcm}{lcm}

\title{Worksheet for Chapter 11}
\author{Egbert Rijke}
\date{October 3, 2025}

\usepackage{hyperref}
\hypersetup{
    colorlinks=true,
    linkcolor=purple,
    filecolor=magenta,
    urlcolor=blue,
    citecolor=purple,
    pdftitle={Lecture notes on Number Theory},
    pdfpagemode=FullScreen,
}
\usepackage[capitalize,english]{cleveref}

  % We set up the exercise environment, which produces list environment in a new
% unnumbered subsection that also gets mentioned in the table of contents.

\newtotcounter{exercisecounter}
\newcommand{\exercise}{\item\stepcounter{exercisecounter}}

\newlist{exenum}{enumerate}{1}
\setlist[exenum]{noitemsep,label=\arabic*.}
  %,ref=\thechapter.\arabic*}

\crefname{exenumi}{Exercise}{Exercises}
  
\newlist{subexenum}{enumerate}{1}
\setlist[subexenum]{noitemsep,label=(\alph*),ref=\theexenumi.\alph*}
\crefname{subexenumi}{Exercise}{Exercises}
  
\newenvironment{exercises}
{%
\subsection*{Exercises}%
\addcontentsline{toc}{section}{Exercises}%
\sectionmark{Exercises}%
\begin{exenum}}
{%
%\addtocounter{exercisecounter}{\theexenumi}%
\end{exenum}}

\newcommand{\floorfrac}[2]{\left\lfloor\frac{#1}{#2}\right\rfloor}

\begin{document}
\maketitle

\thispagestyle{empty}

Recall that for $0\leq k\leq n$ the falling factorial $n_k$ is defined by
\begin{equation*}
  n_k=n(n-1)\cdots (n-k+1).
\end{equation*}
In this worksheet, we will assume that $2k\leq n$. We will also write $m:=n-k$ so that $n=m+k$.

\begin{exercises}
  \exercise Prove that $n_k\nmid k!$.
  \exercise Use Legendre's formula to find an expression for $\nu_p(n_k)$.
  \exercise Consider two nonnegative integers $a=qd+r$ and $b=pd+s$, where $0\leq r,s<d$. Show that
  \begin{equation*}
    \floorfrac{a+b}{d}-\floorfrac{b}{d}-\floorfrac{a\vphantom{b}}{d}=
    \begin{cases}
      0 & \text{if $r+s<d$,}\\
      1 & \text{otherwise.}
    \end{cases}
  \end{equation*}
  \exercise Prove that
    \begin{equation*}
    \left\lfloor\frac{n}{p^j}\right\rfloor-\left\lfloor\frac{m}{p^j}\right\rfloor =
    \begin{cases}
      \floorfrac{k}{p^j} & \text{if $b_j+d_j<p^j$,} \\[1em]
      \floorfrac{k}{p^j} + 1 & \text{if $b_j+d_j\geq p^j$.}
    \end{cases}
  \end{equation*}
  \exercise Prove that for any prime number $p\leq k$, we have
  \begin{equation*}
    {\textstyle \nu_p\left(\binom{n}{k}\right)}\leq 1+\lfloor\log_p(k)\rfloor.
  \end{equation*}
  \exercise Prove that if every prime factor $p\mid n_k$ satisfies $p\leq k$, then we have the inequalities
  \begin{equation*}
    \frac{4^k}{2k+1}<\binom{n}{k}\leq \prod_{p\leq k}p^{1+\lfloor\log_p(k)\rfloor}.
  \end{equation*}
  \exercise (Sylvester--Schur) Conclude that if $2k\leq n$, then $n_k$ has a prime divisor $p>k$.
  \exercise (Bertrand's Postulate) Prove that there exists a prime $n<p<2n$ for every $n$.
  \exercise (Bonse's inequality). Let $p_k$ be the $k$th prime, so that $p_1=2$, $p_2=3$, $p_3=5$, and so forth. Use Bertrand's postulate to show that for $k\geq 5$, then
  \begin{equation*}
    \prod_{k=1}^n p_k < p_{n+1}^3.
  \end{equation*}
\end{exercises}
\end{document}
